\documentclass[border=0pt]{standalone}
\usepackage[RPvoltages]{circuitikz}
\usetikzlibrary{arrows,decorations.markings,arrows.meta}
\newcommand{\bushere}[3]{% length, text above, text below
    coordinate(tmp)
    % go up and do an edge down
    ++(0,#1) node[anchor=base]{#2} edge[line width=2.5pt] ++(0,{-1.88*#1})
    % edges do not move the current point, go down to position the node
    ++(0,{-1.88*#1}) node[below]{#3}
    % go back to where we started
    (tmp)
}

\newcommand{\redbushere}[3]{% length, text above, text below
%    coordinate(tmp)
    % go up and do an edge down
    ++(0.07,#1) node[anchor=base]{#2} edge[red,dashed, line width=1.5pt]  ++(0,{-1.88*#1})
    % edges do not move the current point, go down to position the node
    ++(0,{-1.88*#1}) node[below]{#3}
    % go back to where we started
%    (tmp)
}


\begin{document}\sf
%\ctikzset{sources/fill=gray!20, resistors/fill=gray!20}
%\ctikzset{sources/fill=white, resistors/fill=white}
%\tikzset{myarr/.style={line width=0.74pt,-{Triangle[width=6pt, length=9pt]}}}
\tikzset{myarr/.style={line width=0.74pt,-{Triangle[width=6pt, length=9pt]}}}

\begin{tikzpicture}[>=latex]% default line width
%\draw [myarr](0,0) to (1.5,0);
\coordinate (B) at (1.5,0);
\coordinate (B1) at (1.4,0);
\coordinate (B2) at (1.4,0);
\node[draw, line width=0.76pt, text=white, minimum width=3.6mm, minimum height=17.5mm, xshift=0.4mm] (dos) at (B) {};

\node[xshift=4.3mm, yshift=0.6mm] (e) at (B) {\Large $e$};

\path (B2)   \bushere{1}{}{} ;

\path (B1)   \redbushere{1}{}{} ;

\draw[myarr,red] (B)++(0.22,0.55)node[xshift=5ex,yshift=1.3ex]{$p_d$}to ++(1.4,0);

\draw[myarr] (0,0)node[xshift=4	ex,yshift=1.3ex]{$p_i$}to (B1);

\draw [myarr](B)++(0.22,-0.55)node[xshift=6ex,yshift=1.3ex]{$p_u$}to++(1.6,0) ;


\end{tikzpicture}
\end{document}